\section{Type-Safe Model API}

\label{sec:api}

\subsection{Design Philosophy}

Candle's model definition API leverages Rust's type system to provide compile-time guarantees:

\begin{enumerate}[leftmargin=1.4em]
    \item \textbf{Dimension checking}: Tensor shapes are tracked at the type level where possible, catching shape mismatches at compile time.
    \item \textbf{Device safety}: Operations between tensors on different devices are compile-time errors.
    \item \textbf{Dtype safety}: Operations between incompatible dtypes require explicit casting.
    \item \textbf{Lifetime safety}: Rust's ownership system prevents use-after-free and double-free bugs that plague C++ tensor libraries.
\end{enumerate}

\subsection{Core Traits}

The API is built on three core traits:

\begin{lstlisting}[language=Rust,caption={Core Candle traits.}]
/// A neural network module
pub trait Module {
    fn forward(&self, x: &Tensor)
        -> Result<Tensor>;
}

/// A module with learnable parameters
pub trait ModuleT: Module {
    fn parameters(&self) -> Vec<&Tensor>;
}

/// Quantized module for efficient inference
pub trait QuantizedModule {
    fn forward_q(&self, x: &Tensor)
        -> Result<Tensor>;
    fn quantization(&self) -> Quantization;
}
\end{lstlisting}

\subsection{Model Definition Example}

A transformer layer in Candle:

\begin{lstlisting}[language=Rust,caption={Transformer layer definition.}]
pub struct TransformerBlock {
    attn: MultiHeadAttention,
    ffn: FeedForward,
    norm1: LayerNorm,
    norm2: LayerNorm,
}

impl Module for TransformerBlock {
    fn forward(&self, x: &Tensor)
        -> Result<Tensor> {
        // Pre-norm attention
        let h = self.norm1.forward(x)?;
        let h = self.attn.forward(&h)?;
        let x = (x + h)?;
        // Pre-norm FFN
        let h = self.norm2.forward(&x)?;
        let h = self.ffn.forward(&h)?;
        let x = (x + h)?;
        Ok(x)
    }
}
\end{lstlisting}

\subsection{Error Handling}

Candle uses Rust's \texttt{Result} type for all fallible operations, providing precise error messages:

\begin{itemize}[leftmargin=1.1em]
    \item \texttt{ShapeMismatch \{expected: [B, S, D], got: [B, D]\}}: Dimension error with full context.
    \item \texttt{DtypeMismatch \{op: "matmul", lhs: f16, rhs: f32\}}: Type mismatch with operation context.
    \item \texttt{DeviceMismatch \{lhs: Cuda(0), rhs: Cpu\}}: Device placement error.
    \item \texttt{OutOfMemory \{requested: 4GB, available: 2.1GB, device: Cuda(0)\}}: Memory error with allocation context.
\end{itemize}

Errors propagate via the \texttt{?} operator with zero runtime cost on the success path (Rust's zero-cost error handling).

\subsection{Compile-Time Guarantees}

\begin{theorem}[Type Safety]
A Candle program that compiles without \texttt{unsafe} blocks satisfies the following invariants:
\begin{enumerate}[leftmargin=1.4em]
    \item No use-after-free: Tensor storage is valid for the lifetime of any reference.
    \item No data races: Concurrent access requires \texttt{Arc<Mutex<T>>} or \texttt{Arc<RwLock<T>>}.
    \item No null pointers: All tensor references are guaranteed non-null by Rust's type system.
    \item No buffer overflows: Array bounds are checked at runtime (and compile-time where shapes are static).
\end{enumerate}
\end{theorem}

\begin{proof}
These invariants follow directly from Rust's ownership, borrowing, and lifetime rules as proven by RustBelt~\cite{jung2018rustbelt}.
\end{proof}

